% TOP_PROBLEM: {"id": 2142, "title": "Erdős Problem #488", "category_id": 1, "category": {"id": 1, "name": "number_theory", "display_name": "Number Theory", "description": "Properties of integers, prime numbers, Diophantine equations.", "slug": "number-theory", "order_index": 1, "created_at": "2026-07-31T15:26:25.670Z"}, "status": "open"}
% FIRST_POSTED: null
% ENTRY_KIND: statement_only
% Imported from: unsolvedmath_additional_500.tex; UnsolvedMath ID 2142
% !TeX program = lualatex
% Compile twice: lualatex 2142.tex
% Self-contained: no external bibliography, images, or \input files.
% Numeric section labels and visible numbers are original UnsolvedMath IDs.
\documentclass[11pt,a4paper]{article}
\usepackage[margin=24mm,headheight=14pt]{geometry}
\usepackage{fontspec}
\setmainfont{Latin Modern Roman}
\setsansfont{Latin Modern Sans}
\setmonofont{Latin Modern Mono}
\usepackage{amsmath,amssymb,mathtools}
\usepackage{unicode-math}
\setmathfont{Latin Modern Math}
\usepackage{microtype}
\usepackage{enumitem,xurl,needspace}
\usepackage[dvipsnames]{xcolor}
\usepackage{hyperref}
\hypersetup{unicode=true,colorlinks=true,linkcolor=MidnightBlue,urlcolor=MidnightBlue,
 pdftitle={UnsolvedMath Problem 2142},pdfauthor={Source-annotated compilation},bookmarksnumbered=true}
\usepackage{bookmark,tocloft,titlesec,fancyhdr}
\setlength{\cftsecnumwidth}{4.2em}
\cftsetpnumwidth{2.5em}
\cftsetrmarg{4em}
\titleformat{\section}{\Large\bfseries\raggedright}{\thesection}{0.7em}{}
\pagestyle{fancy}\fancyhf{}
\fancyhead[L]{\small\sffamily Additional UnsolvedMath statements}
\fancyhead[R]{\small Original catalogue IDs}
\fancyfoot[C]{\thepage}
\setlength{\parindent}{0pt}
\setlength{\parskip}{5pt plus 1pt}
\setlength{\emergencystretch}{3em}
\setcounter{tocdepth}{1}\setcounter{secnumdepth}{1}
\setlist[itemize]{leftmargin=3em,itemsep=4pt,topsep=4pt,parsep=0pt}
\allowdisplaybreaks
\newcommand{\N}{\mathbb{N}}
\newcommand{\Z}{\mathbb{Z}}
\newcommand{\Q}{\mathbb{Q}}
\newcommand{\R}{\mathbb{R}}
\newcommand{\C}{\mathbb{C}}
\newcommand{\F}{\mathbb{F}}
\newcommand{\A}{\mathbb{A}}
\newcommand{\PP}{\mathbb{P}}
\newcommand{\E}{\mathbb{E}}
\newcommand{\eps}{\varepsilon}
\newcommand{\dd}{\,\mathrm d}
\newcommand{\sref}[2]{\hyperlink{source:#1:#2}{[S#2]}}
\newif\ifshowcatalogue
\showcataloguetrue
\newcommand{\cataloguescope}[1]{\ifshowcatalogue
 \paragraph{Statement recorded in the supplied source.}
 #1\fi}
\begin{document}
% UnsolvedMath ID 2142; source catalogue code EP-488

\setcounter{section}{2141}

\section{A Factor-Two Bound for Prefix Densities of Multiples}\label{2142}

{\small\sffamily UnsolvedMath ID 2142 · Number Theory}

\subsection{Definitions and mathematical statement}

\paragraph{Shared notation.}Unless a section says otherwise, $\N=\{1,2,\ldots\}$,
$\Z$ is the ring of integers, $\Q,\R,\C$ are the rational, real, and complex
fields, $[N]=\{1,\ldots,\lfloor N\rfloor\}$, and $\log$ is the natural
logarithm. For real functions of a parameter tending to infinity,
$f\ll g$ and $f=O(g)$ mean $|f|\leq Cg$ eventually for some fixed $C$;
$f\gg g$ reverses this comparison; $f=o(g)$ means $f/g\to0$;
$f\sim g$ means $f/g\to1$; and $f\asymp g$ means both order bounds.
A subscript on $O$ or $\ll$ specifies allowed constant dependence.
The expression $n^{o(1)}$ in an upper bound means growth smaller than
$n^\epsilon$ for each fixed $\epsilon>0$. Local definitions govern any
exception, finite-field convention, or use of the same symbol for another
object. An asymptotic claim is not automatically an effective algorithm.



Assume $A$ is a nonempty finite set of positive integers so that $\max A$ exists. The multiples set is periodic, but the requested inequality compares all finite-prefix densities beyond the largest specified divisor, not just their eventual limit. \sref{2142}{1} \sref{2142}{2}

\paragraph{Statement recorded in the supplied source.}
Let $A$ be a finite set and $ B=\{ n \geq 1 : a\mid n\textrm{ for some }a\in A\}. $ Is it true that, for every $m>n\geq \max(A)$, $ \frac{\lvert B\cap [1,m]\rvert }{m}< 2\frac{\lvert B\cap [1,n]\rvert}{n}? $ \sref{2142}{1}

{\small\emph{Transcription note.} Only unambiguous escape/formatting damage and nonmathematical serialization debris were repaired in the displayed source statement.}

\subsection{Short English statement}

Once an interval passes all prescribed divisors, can the proportion of integers divisible by at least one of them ever double at a later endpoint?

\subsection{Sources}

{\small

\begin{itemize}

\item[\hypertarget{source:2142:1}{S1}] ULAM.AI, \emph{UnsolvedMath}, supplied \texttt{problems.json}, record 2142 (EP-488), statement and background. The embedded literature-review date is 2026-08-17; that is the archive’s review date, not a fresh certification. Dataset landing page: \url{https://huggingface.co/datasets/ulamai/UnsolvedMath}.

\item[\hypertarget{source:2142:2}{S2}] T. F. Bloom, \emph{Erdős Problems}, Problem 488, \url{https://www.erdosproblems.com/488}. Reference imported from the supplied record; the live problem page was not independently checked for this compilation.

\item[\hypertarget{source:2142:3}{S3}] P. Erdős, Some Unsolved Problems, Magyar Tud. Akad. Mat. Kutató Int. Közl. 6 (1961), 221--254. \url{https://real.mtak.hu/201135/}. Bibliographic information imported from the supplied record.

\item[\hypertarget{source:2142:4}{S4}] HUN-REN Rényi Institute bibliography entry for P. Erdős, Remarks on number theory V. Extremal problems in number theory II, Mat. Lapok 17 (1966), 135--155. (bibliographic\_index). \url{https://renyi.hu/en/node/4094}. Bibliographic information imported from the supplied record.

\item[\hypertarget{source:2142:5}{S5}] Erd\H{o}s, P\'al, Remarks on number theory. {V}. {E}xtremal problems in number theory. {II}. Mat. Lapok (1966), 135--155. Bibliographic information imported from the supplied record.

\item[\hypertarget{source:2142:6}{S6}] Guy, Richard K., Unsolved problems in number theory. (2004), xviii+437. Bibliographic information imported from the supplied record.

\end{itemize}

}

\label{end:2142}
\end{document}
